%! TEX program = lualatex
\documentclass[10pt, a4paper]{beamer} %handout fuer eine gedruckte version der slides
\usetheme{metropolis}

\usepackage{url}
\usepackage{polyglossia}
\setmainlanguage{english}

\usepackage{listings,xcolor}
\usepackage{graphicx}
\usepackage{nameref}
\usepackage{mathtools}
\usepackage{tikz}

\usepackage{booktabs}

\makeatletter
\newcommand*{\currentname}{\@currentlabelname}
\makeatother

\usepackage{fontspec}
\newfontfamily\Bera{Fira Code}[Scale=0.85]

\setmonofont[
  Contextuals={Alternate}
]{Fira Code}


\definecolor{pdfToolsBlue}{HTML}{0084C7}
\definecolor{pdfToolsDarkBlue}{HTML}{004569}
\definecolor{pdfToolsDarkerBlue}{HTML}{002A3F}
\definecolor{pdfToolsRenderingOrange}{HTML}{E78B00}

\definecolor{mDarkTeal}{HTML}{23373b}
\definecolor{mLightBrown}{HTML}{EB811B}
\definecolor{mDarkBlue}{HTML}{0F3470}
\definecolor{mDarkGrey}{HTML}{999999}
\definecolor{mLighterGrey}{HTML}{CCCCCC}
\definecolor{mLightGrey}{rgb}{0.95,0.95,0.95}
\definecolor{mLightBlue}{HTML}{212751}

\newcommand{\lb}[1]{{\color{pdfToolsRenderingOrange}#1}}

\newcommand{\bblock}[3]{
        \begin{block}{#1}
        #2
    \end{block}
}

\setbeamercolor{block title}{fg=pdfToolsDarkBlue}
\setbeamercolor{frametitle}{bg=pdfToolsBlue,fg=white}
\setbeamercolor{normal text}{fg=black}

\lstset{language=C++,
    backgroundcolor=\color{mLightGrey},
    frame=single,
    rulecolor=\color{mLightGrey}, % make frame "invisible"
    basicstyle=\Bera\tiny,
    keywordstyle=\bfseries\color{mDarkBlue},
    commentstyle=\color{mDarkGrey}\itshape,
    captionpos=t,
    emphstyle=\ttb\color{mDarkBlue},    % Custom highlighting style
    stringstyle=\color{mLightBrown},
    tabsize=2,
    numberbychapter=false,
    showstringspaces=false,
    breaklines=true,
    morekeywords={assert, break, class, continue, else, except, for, if, return, try, while}
}
%
\makeatletter
\setbeamertemplate{footline}{%
\leavevmode
\vbox{\begin{beamercolorbox}[dp=1.25ex,ht=2.75ex]{fg=black}%
  \hspace*{1em}\insertsectionhead%
  \ifx\insertsubsectionhead\@empty\relax\else\hspace*{\fill}\insertsubsectionhead\hspace*{2ex}\fi
  \end{beamercolorbox}%
  }%
}
\makeatother


\newcommand{\itembullet}{
    \tikz\draw[pdfToolsDarkerBlue,fill=pdfToolsDarkerBlue] (0,0) circle (.4ex);
}

\AtBeginSection{}

\setbeamertemplate{section in toc}{%
  {\color{mDarkTeal}\rule[0.3ex]{3pt}{3pt}}~\inserttocsection\par}
\setbeamertemplate{subsection in toc}{%
  \hspace{1.2em}{\color{mLightBrown}\rule[0.3ex]{3pt}{3pt}}~\inserttocsubsection\par}


\setbeamertemplate{itemize item}{\itembullet}
\setlength{\parskip}{4ex}

\title % (optional, only for long titles)
{Dynamic Loading of Modules in WASM/Emscripten}
\author % (optional, for multiple authors)
{Philipp Gloor}
\date{}
\titlegraphic{\raggedleft\includegraphics[width=0.2\textwidth]{img/logo-pdftools-c-200}}
\begin{document}
\begin{frame}
\titlepage
\end{frame}

{
  \setlength{\parskip}{0ex}
\begin{frame}{Table of Contents}
    \tableofcontents
  \end{frame}
}
\section{Emscripten: Main Module, Side Modules}
\begin{frame}[c, allowframebreaks]\frametitle{Emscripten: Main Module, Side Modules}
  In order to be able to dynamically load a WASM file they need to be compiled with 
  specific options.
  
  Following options can control how Emscripten creates the WebAssembly module:
  \begin{itemize}
    \item \texttt{-s MAIN\_MODULE}
    \item \texttt{-s SIDE\_MODULE}
  \end{itemize}

  The differences between the main module and the side modules:
\framebreak
  \begin{itemize}
    \item Main modules have the system libraries linked in
    \item Side modules are pure wasm files that contain LLVM bitcode that have no system libraries
    \item Compiling a main module adds a JavaScript file which sets up the Emscripten environment
  \end{itemize}
\end{frame}

\section{Dynamic Loading}
\begin{frame}[allowframebreaks]{Dynamic Loading}
  Three different methods to dynamically link functions:

  \begin{itemize}
    \item dlopen/dlsym: Works the same as with native C code. Requires the side module to be available in the file system. Relies on C-Linkage or using mangled
    C++ names. Obtain dll handle with \texttt{dlopen} and link the functions with \texttt{dlsym}
    \item On startup: Define \texttt{Module.dynamicLibraries = [<list of wasm names>]}. Loads the wasm modules synchronously. Functions from modules need 
    to be declared for compilation to work.
    \item On demand: Use \texttt{EM\_ASM()} macro within C/C++.
  \end{itemize}
  \framebreak

    Using \texttt{loadDynamicLibrary} from within C/C++ seems the best way

    \begin{enumerate}
      \item Loads WASM when really needed
      \item Automatic linking (unlike dlopen)
      \item Asynchronous
    \end{enumerate}
\end{frame}

\section{Passing values: Reference or Value}
\begin{frame}{Passing values: Reference or Value}

\end{frame}

\section{}
\end{document}